%%%%%%%%%%%%%%%%
%%% gen 5 %%%
%%%%%%%%%%%%%%%%

\Chapter{5}

\Verse{1}
{
	\hBookOfRecords{זֶה סֵפֶר תּוֹלְדֹת אָדָם בְּיוֹם בְּרֹא אֱלֹהִים אָדָם בִּדְמוּת אֱלֹהִים עָשָׂה אֹתוֹ׃}
}
{
	\eBookOfRecords{
		5 This is the Book of Records of the Human. In the day of
		{\God}’s creating a human, He made it in the likeness of {\God}.
	}
}
{
%	\footnote{
%		Records. This word has usually been understood to mean “generations,”
%		but that is inadequate. The word is used both to introduce records of
%		births (as in this verse) and to introduce stories of events within a
%		family (as in Gen 37:2). Footnote 5:1. human. In these verses, the
%		Hebrew ’dm sometimes refers to the species, as in v. 2, where it applies
%		to both male and female; and it sometimes refers specifically to the
%		first male human, so that it also functions as his personal name: Adam.
%	}
}

\Verse{2}
{
	\hBookOfRecords{זָכָר וּנְקֵבָה בְּרָאָם וַיְבָרֶךְ אֹתָם וַיִּקְרָא אֶת שְׁמָם אָדָם בְּיוֹם הִבָּרְאָם׃}
}
{
	\eBookOfRecords{
		He created them male and female, and He blessed them and
		called their name “Human” in the day of their being
		created.
	}
}
{
}

\Verse{3}
{
	\hBookOfRecords{וַיְחִי אָדָם שְׁלֹשִׁים וּמְאַת שָׁנָה וַיּוֹלֶד בִּדְמוּתוֹ כְּצַלְמוֹ וַיִּקְרָא אֶת שְׁמוֹ שֵׁת׃}
}
{
	\eBookOfRecords{
		And the human lived a hundred thirty years, and he
		fathered in his likeness—like his image—and called his
		name Seth.
	}
}
{
%	\footnote{
%		he fathered in his likeness—like his image. The first man’s similarity
%		to his son is described with the same two nouns that are used to
%		describe the first two humans’ similarity to God (1:26–27). It certainly
%		sounds as if it means something physical here. We surely would have
%		taken it that way if we had read this verse without having read Genesis
%		1. Still, we must be cautious on such a classic biblical question. In
%		any case, the significance of this verse is to establish that whatever
%		it is that the first humans acquire from God, it is something that
%		passes by heredity. It is not only the first two humans but the entire
%		species that bears God’s image. Footnote 5:3. Seth. With the birth of
%		Adam’s and Eve’s son Seth, the text begins a flow of generations,
%		narrowing to a particular family, that continues unbroken through the
%		book of Genesis and ultimately through the rest of the Hebrew Bible.
%		Ironically, the element that establishes this flow, that produces the
%		continuity of Genesis, and that sets the history of the family into the
%		context of the universal history is the “begat” lists. It is ironic
%		because these lists are tedious to most readers. They therefore skip,
%		skim, plow through, or joke about them. The result is that many (perhaps
%		most) readers never get the feeling of Genesis as a book. It is a
%		continuous, sensible work, filled with connections, ironies, puns, and
%		character development—which are diminished or even lost when one reads
%		it only as a collection of separate stories.
%	}
}

\Verse{4}
{
	\hBookOfRecords{וַיִּהְיוּ יְמֵי אָדָם אַחֲרֵי הוֹלִידוֹ אֶת שֵׁת שְׁמֹנֶה מֵאֹת שָׁנָה וַיּוֹלֶד בָּנִים וּבָנוֹת׃}
}
{
	\eBookOfRecords{
		And the human’s days after his fathering Seth were eight
		hundred years, and he fathered sons and daughters.
	}
}
{
%	\footnote{
%		he fathered sons and daughters. This is the presumed answer to the
%		question of where Cain’s wife came from.
%	}
}

\Verse{5}
{
	\hBookOfRecords{וַיִּהְיוּ כּל יְמֵי אָדָם אֲשֶׁר חַי תְּשַׁע מֵאוֹת שָׁנָה וּשְׁלֹשִׁים שָׁנָה וַיָּמֹת׃}
}
{
	\eBookOfRecords{
		And all of the human’s days that he lived were nine
		hundred years and thirty years. And he died.
	}
}
{
%	\footnote{
%		nine hundred years and thirty years. The long life spans in the early
%		portions of the Torah are an old question. Some assume that the ancients
%		must have counted years differently. But that is not correct. (If we
%		divide Adam’s 930 years by ten to get it within normal range, then how
%		shall we divide Moses’ 120?) It is clear that this author thought of a
%		year as a normal solar year because that is how long the flood lasts.
%		The point to note is: life spans are pictured as growing shorter. The
%		ten generations from Adam to Noah approach ages of 1,000. But the last
%		one to live more than 900 years is Noah. The next ten generations start
%		with Shem, who lives 600 years, and life spans decline after him. The
%		last person to live more than 200 years is Terah. Abraham (175), Isaac
%		(180), and Jacob (147) live long lives, but not as long as their
%		ancestors. And Moses lives to be 120, which is understood to have
%		become, at some point, the maximum for human life. (See the comment on
%		Gen 6:3.)
%	}
}

\Verse{6}
{
	\hBookOfRecords{וַיְחִי שֵׁת חָמֵשׁ שָׁנִים וּמְאַת שָׁנָה וַיּוֹלֶד אֶת אֱנוֹשׁ׃}
}
{
	\eBookOfRecords{
		And Seth lived five years and a hundred years, and he
		fathered Enosh.
	}
}
{
}

\Verse{7}
{
	\hBookOfRecords{וַיְחִי שֵׁת אַחֲרֵי הוֹלִידוֹ אֶת אֱנוֹשׁ שֶׁבַע שָׁנִים וּשְׁמֹנֶה מֵאוֹת שָׁנָה וַיּוֹלֶד בָּנִים וּבָנוֹת׃}
}
{
	\eBookOfRecords{
		And Seth lived after his fathering Enosh seven years and
		eight hundred years, and he fathered sons and daughters.
	}
}
{
}

\Verse{8}
{
	\hBookOfRecords{וַיִּהְיוּ כּל יְמֵי שֵׁת שְׁתֵּים עֶשְׂרֵה שָׁנָה וּתְשַׁע מֵאוֹת שָׁנָה וַיָּמֹת׃}
}
{
	\eBookOfRecords{
		And all of Seth’s days were twelve years and nine hundred
		years. And he died.
	}
}
{
}

\Verse{9}
{
	\hBookOfRecords{וַיְחִי אֱנוֹשׁ תִּשְׁעִים שָׁנָה וַיּוֹלֶד אֶת קֵינָן׃}
}
{
	\eBookOfRecords{And Enosh lived ninety years, and he fathered Cainan.}
}
{
}

\Verse{10}
{
	\hBookOfRecords{וַיְחִי אֱנוֹשׁ אַחֲרֵי הוֹלִידוֹ אֶת קֵינָן חֲמֵשׁ עֶשְׂרֵה שָׁנָה וּשְׁמֹנֶה מֵאוֹת שָׁנָה וַיּוֹלֶד בָּנִים וּבָנוֹת׃}
}
{
	\eBookOfRecords{
		And Enosh lived after his fathering Cainan fifteen years
		and eight hundred years, and he fathered sons and
		daughters.
	}
}
{
}

\Verse{11}
{
	\hBookOfRecords{וַיִּהְיוּ כּל יְמֵי אֱנוֹשׁ חָמֵשׁ שָׁנִים וּתְשַׁע מֵאוֹת שָׁנָה וַיָּמֹת׃}
}
{
	\eBookOfRecords{
		And all of Enosh’s days were five years and nine hundred
		years. And he died.
	}
}
{
}

\Verse{12}
{
	\hBookOfRecords{וַיְחִי קֵינָן שִׁבְעִים שָׁנָה וַיּוֹלֶד אֶת מַהֲלַלְאֵל׃}
}
{
	\eBookOfRecords{And Cainan lived seventy years, and he fathered Mahalalel.}
}
{
}

\Verse{13}
{
	\hBookOfRecords{וַיְחִי קֵינָן אַחֲרֵי הוֹלִידוֹ אֶת מַהֲלַלְאֵל אַרְבָּעִים שָׁנָה וּשְׁמֹנֶה מֵאוֹת שָׁנָה וַיּוֹלֶד בָּנִים וּבָנוֹת׃}
}
{
	\eBookOfRecords{
		And Cainan lived after his fathering Mahalalel forty years
		and eight hundred years, and he fathered sons and
		daughters.
	}
}
{
}

\Verse{14}
{
	\hBookOfRecords{וַיִּהְיוּ כּל יְמֵי קֵינָן עֶשֶׂר שָׁנִים וּתְשַׁע מֵאוֹת שָׁנָה וַיָּמֹת׃}
}
{
	\eBookOfRecords{
		And all of Cainan’s days were ten years and nine hundred
		years. And he died.
	}
}
{
}

\Verse{15}
{
	\hBookOfRecords{וַיְחִי מַהֲלַלְאֵל חָמֵשׁ שָׁנִים וְשִׁשִּׁים שָׁנָה וַיּוֹלֶד אֶת יָרֶד׃}
}
{
	\eBookOfRecords{
		And Mahalalel lived five years and sixty years, and he
		fathered Jared.
	}
}
{
}

\Verse{16}
{
	\hBookOfRecords{וַיְחִי מַהֲלַלְאֵל אַחֲרֵי הוֹלִידוֹ אֶת יֶרֶד שְׁלֹשִׁים שָׁנָה וּשְׁמֹנֶה מֵאוֹת שָׁנָה וַיּוֹלֶד בָּנִים וּבָנוֹת׃}
}
{
	\eBookOfRecords{
		And Mahalalel lived after his fathering Jared thirty years
		and eight hundred years, and he fathered sons and
		daughters.
	}
}
{
}

\Verse{17}
{
	\hBookOfRecords{וַיִּהְיוּ כּל יְמֵי מַהֲלַלְאֵל חָמֵשׁ וְתִשְׁעִים שָׁנָה וּשְׁמֹנֶה מֵאוֹת שָׁנָה וַיָּמֹת׃}
}
{
	\eBookOfRecords{
		And all of Mahalalel’s days were ninety-five years and
		eight hundred years. And he died.
	}
}
{
}

\Verse{18}
{
	\hBookOfRecords{וַיְחִי יֶרֶד שְׁתַּיִם וְשִׁשִּׁים שָׁנָה וּמְאַת שָׁנָה וַיּוֹלֶד אֶת חֲנוֹךְ׃}
}
{
	\eBookOfRecords{
		And Jared lived sixty-two years and a hundred years, and
		he fathered Enoch.
	}
}
{
}

\Verse{19}
{
	\hBookOfRecords{וַיְחִי יֶרֶד אַחֲרֵי הוֹלִידוֹ אֶת חֲנוֹךְ שְׁמֹנֶה מֵאוֹת שָׁנָה וַיּוֹלֶד בָּנִים וּבָנוֹת׃}
}
{
	\eBookOfRecords{
		And Jared lived after his fathering Enoch eight hundred
		years, and he fathered sons and daughters.
	}
}
{
}

\Verse{20}
{
	\hBookOfRecords{וַיִּהְיוּ כּל יְמֵי יֶרֶד שְׁתַּיִם וְשִׁשִּׁים שָׁנָה וּתְשַׁע מֵאוֹת שָׁנָה וַיָּמֹת׃}
}
{
	\eBookOfRecords{
		And all of Jared’s days were sixty-two years and nine
		hundred years. And he died.
	}
}
{
}

\Verse{21}
{
	\hBookOfRecords{וַיְחִי חֲנוֹךְ חָמֵשׁ וְשִׁשִּׁים שָׁנָה וַיּוֹלֶד אֶת מְתוּשָׁלַח׃}
}
{
	\eBookOfRecords{
		And Enoch lived sixty-five years, and he fathered
		Methuselah.
	}
}
{
}

\Verse{22}
{
	\hBookOfRecords{וַיִּתְהַלֵּךְ חֲנוֹךְ אֶת הָאֱלֹהִים אַחֲרֵי הוֹלִידוֹ אֶת מְתוּשֶׁלַח שְׁלֹשׁ מֵאוֹת שָׁנָה וַיּוֹלֶד בָּנִים וּבָנוֹת׃}
}
{
	\eBookOfRecords{
		And Enoch walked with {\God} after his fathering Methuselah
		three hundred years, and he fathered sons and daughters.
	}
}
{
}

\Verse{23}
{
	\hBookOfRecords{וַיְהִי כּל יְמֵי חֲנוֹךְ חָמֵשׁ וְשִׁשִּׁים שָׁנָה וּשְׁלֹשׁ מֵאוֹת שָׁנָה׃}
}
{
	\eBookOfRecords{
		And all of Enoch’s days were sixty-five years and three
		hundred years.
	}
}
{
}

\Verse{24}
{
	\hBookOfRecords{וַיִּתְהַלֵּךְ חֲנוֹךְ אֶת הָאֱלֹהִים וְאֵינֶנּוּ כִּי לָקַח אֹתוֹ אֱלֹהִים׃}
}
{
	\eBookOfRecords{
		And Enoch walked with {\God}, and he was not, because {\God}
		took him.
	}
}
{
%	\footnote{
%		Enoch walked with God. This expression is used in ancient Near Eastern
%		texts to express continuous fidelity. So here it would mean that Enoch
%		was faithful to God. Footnote 5:24. and he was not, because God took
%		him. I do not know what this means. It was traditionally understood to
%		mean that Enoch does not die. Alternatively, it could be the report of
%		his death. It comes at the point at which all the other entries in this
%		list say “And he died.” The same word is used later by Joseph’s brothers
%		to express the fact that he is gone (Gen 42:13). At that point, the
%		brothers do not know whether he is alive or not. It may possibly mean
%		something like that in the case of Enoch as well: his fate was unknown.
%		At minimum, it means (as with the case of Elijah later) that there is
%		something distinctive and unusual about the man, his relationship to
%		God, and his departure from life among humans.
%	}
}

\Verse{25}
{
	\hBookOfRecords{וַיְחִי מְתוּשֶׁלַח שֶׁבַע וּשְׁמֹנִים שָׁנָה וּמְאַת שָׁנָה וַיּוֹלֶד אֶת לָמֶךְ׃}
}
{
	\eBookOfRecords{
		And Methuselah lived eighty-seven years and a hundred
		years, and he fathered Lamech.
	}
}
{
}

\Verse{26}
{
	\hBookOfRecords{וַיְחִי מְתוּשֶׁלַח אַחֲרֵי הוֹלִידוֹ אֶת לֶמֶךְ שְׁתַּיִם וּשְׁמוֹנִים שָׁנָה וּשְׁבַע מֵאוֹת שָׁנָה וַיּוֹלֶד בָּנִים וּבָנוֹת׃}
}
{
	\eBookOfRecords{
		And Methuselah lived after his fathering Lamech eighty-two
		years and seven hundred years, and he fathered sons and
		daughters.
	}
}
{
}

\Verse{27}
{
	\hBookOfRecords{וַיִּהְיוּ כּל יְמֵי מְתוּשֶׁלַח תֵּשַׁע וְשִׁשִּׁים שָׁנָה וּתְשַׁע מֵאוֹת שָׁנָה וַיָּמֹת׃}
}
{
	\eBookOfRecords{
		And all of Methuselah’s days were sixty-nine years and
		nine hundred years. And he died.
	}
}
{
}

\Verse{28}
{
	\hBookOfRecords{וַיְחִי לֶמֶךְ שְׁתַּיִם וּשְׁמֹנִים שָׁנָה וּמְאַת שָׁנָה וַיּוֹלֶד בֵּן׃}
}
{
	\eBookOfRecords{
		And Lamech lived eighty-two years and a hundred years, and
		he fathered a son
	}
}
{
}

\Verse{29}
{
	\hR{וַיִּקְרָא אֶת שְׁמוֹ נֹחַ לֵאמֹר זֶה יְנַחֲמֵנוּ מִמַּעֲשֵׂנוּ וּמֵעִצְּבוֹן יָדֵינוּ מִן הָאֲדָמָה אֲשֶׁר אֵרְרָהּ יְהֹוָה׃}
}
{
	\eR{
		and called his name Noah, saying, “This one will console
		us from our labor and from our hands’ suffering from the
		ground, which YHWH has cursed.”
	}
}
{
%	\footnote{
%		Noah. The name is connected here to the Hebrew root nm, meaning
%		“console,” though we would naturally connect it to a different root, nw,
%		meaning “rest,” which matches the name Noah and does not have the extra
%		m (Hebrew mem) at the end. Biblical names, like contemporary naming of
%		Jewish children, are not necessarily based on precise etymologies, but
%		rather may be based on similarity of sounds, involving only some of the
%		root letters. See, similarly, the comment on the name Cain (Gen 4:1).
%	}
}

\Verse{30}
{
	\hBookOfRecords{וַיְחִי לֶמֶךְ אַחֲרֵי הוֹלִידוֹ אֶת נֹחַ חָמֵשׁ וְתִשְׁעִים שָׁנָה וַחֲמֵשׁ מֵאֹת שָׁנָה וַיּוֹלֶד בָּנִים וּבָנוֹת׃}
}
{
	\eBookOfRecords{
		And Lamech lived after his fathering Noah ninety-five
		years and five hundred years, and he fathered sons and
		daughters.
	}
}
{
}

\Verse{31}
{
	\hBookOfRecords{וַיְהִי כּל יְמֵי לֶמֶךְ שֶׁבַע וְשִׁבְעִים שָׁנָה וּשְׁבַע מֵאוֹת שָׁנָה וַיָּמֹת׃}
}
{
	\eBookOfRecords{
		And all of Lamech’s days were seventy-seven years and
		seven hundred years. And he died.
	}
}
{
}

\Verse{32}
{
	\hBookOfRecords{וַיְהִי נֹחַ בֶּן חֲמֵשׁ מֵאוֹת שָׁנָה וַיּוֹלֶד נֹחַ אֶת שֵׁם אֶת חָם וְאֶת יָפֶת׃}
}
{
	\eBookOfRecords{
		And Noah was five hundred years old, and Noah fathered
		Shem, Ham, and Yaphet.
	}
}
{
}
